\documentclass[../main]{subfiles}
\begin{document}
\ifSubfilesClassLoaded{\setcounter{page}{4}}{}
\section{Il comunismo come realtà nella produzione sociale}
\label{sec:04_comunismo_realta}

L’affermazione secondo cui il comunismo sia già una realtà potrebbe apparire audace. In ogni caso, merita un’attenta verifica. A prima vista, l’evidenza suggerisce il contrario: la sua assenza, l’onnipresenza del lavoro, il cui valore per l’individuo si riduce unicamente alla retribuzione, al salario. Si potrebbe concludere che la cellula – l’embrione – del comunismo, ovvero la forma di attività umana libera e collettiva, rappresenti un elemento marginale nel percorso evolutivo della società umana. Tuttavia, cercheremo di dimostrare che tale conclusione è errata. A tal fine, non basta limitarsi a individuare esempi che attestino la presenza del comunismo in forme embrionali all’interno di una società di classe basata sullo sfruttamento. È necessario analizzare la forma di attività libera e collettiva che si sviluppa nella società di classe contemporanea, valutandone le potenzialità di ulteriore espansione.

È bene ricordare come, in passato, Lenin abbia saputo valutare positivamente, innanzitutto, l'esperienza dei soviet, e successivamente, l'istituzione delle giornate di lavoro comuniste. «In tale contesto, riveste un'importanza capitale l'organizzazione, per iniziativa dei lavoratori stessi, delle giornate di lavoro comuniste. Pur sembrando soltanto un inizio, esso rappresenta in realtà un passo di straordinaria rilevanza. È l'avvio di una rivoluzione ben più ardua, profonda, radicale e decisiva della semplice caduta della borghesia, poiché essa si configura come una vittoria sull'inerzia, sulla passività, sull'egoismo residuo, sulle abitudini che il capitalismo, fonte di ogni male, ha lasciato in eredità al proletariato e al mondo contadino.Solo quando questa vittoria sarà saldamente acquisita, allora – e solo allora – sarà possibile dare vita a una nuova forma di organizzazione sociale, a una disciplina socialista; allora, e solo allora, un ritorno al capitalismo diverrà impraticabile e il comunismo si affermerà in maniera inequivocabile. Lenin riteneva che la prima giornata di lavoro volontario, promossa il 10 maggio 1919 dai ferrovieri della linea Mosca-Kazan, rivestisse un’importanza storica superiore a qualsiasi successo ottenuto durante la Prima Guerra Mondiale. Non si trattava semplicemente di un atto di volontà autonomo, compiuto dai lavoratori che avevano organizzato tale giornata, ma anche di un significativo incremento della produttività, derivante dal fatto che ciascun lavoratore aveva preso a cuore l'iniziativa, vedendo in essa la possibilità di contribuire, con un impegno personale, allo sviluppo della rete ferroviaria, a beneficio dell'intera società.

In seguito alla pubblicazione del capitolo dedicato alle origini del comunismo, uno dei nostri lettori ha sollevato una questione di notevole importanza: «L'organizzazione produttiva comunista implica, in ultima analisi, anche la creazione di beni materiali, e non si limita a un mero esercizio letterario alla maniera di Puškin, finalizzato alla coltivazione di talenti poetici. Oggi, infatti, assistiamo alla proliferazione di presunte "comunità comuniste" fondate sullo sfruttamento a beneficio di una ristretta élite», argomentava il lettore. Chiedeva, pertanto, di fornire delle evidenze concrete dell'esistenza del comunismo e di illustrare come questo potrebbe concretizzarsi, in particolare, nell'ambito della produzione materiale. A nostro avviso, la ragione di tale osservazione risiede in due elementi. In primo luogo, il lettore ha espresso un punto di vista che riflette le istanze del proletariato, ovvero di coloro che sono direttamente coinvolti nella produzione materiale attraverso la vendita della propria forza lavoro, in un contesto in cui i principi del comunismo non trovano ancora piena applicazione.In secondo luogo, il lettore pone un interrogativo fondamentale sulla produzione materiale nell'ambito del comunismo, poiché è proprio questa che ne definisce l'essenza a livello sociale, e non gli ideali, le aspirazioni o perfino le azioni dei rivoluzionari.

Il comunismo, tuttavia, risulta impraticabile in un sistema produttivo esclusivamente orientato alla materialità. In un’epoca in cui l’uomo si dedica primariamente alla fabbricazione di beni, e la genesi delle relazioni sociali avviene in maniera del tutto estranea alle sue intenzioni, ai suoi scopi, alla sua volontà e alla sua coscienza, parlare di comunismo appare utopistico, pur riconoscendo che il processo produttivo possa, in sé, procurare al lavoratore una certa soddisfazione. In questo scenario, l’uomo si riduce a una mera forza materiale, equiparabile ad altre forze naturali. In ultima analisi, egli si comporta come un oggetto che agisce su altri oggetti. Da una prospettiva marxista, in vista dell’avvento del comunismo, questo genere di attività lavorativa non solo può, ma deve essere soppresso in quanto tale, e interamente affidato alle macchine.Analizzando il ruolo storico del capitale, Marx evidenziò in particolare l'importanza della creazione del plusvalore e dell'universalità dei bisogni umani come condizione essenziale e presupposto del comunismo. Il capitale, anche all'interno del sistema capitalista, si occupa già di questo aspetto. Tuttavia, ciò non è sufficiente. Il ruolo storico del capitale si compirà appieno e la società comunista diventerà realtà solo quando verrà superato quel tipo di lavoro in cui l'uomo si limita a svolgere compiti che potrebbero essere eseguiti dalle macchine, agendo quindi per il bene dell'umanità. In altre parole, l'uomo potrà agire in modo veramente comunista solo quando sarà liberato dalla mera produzione materiale, quando non si limiterà a fare ciò che potrebbero fare le macchine e, soprattutto, quando smetterà di comportarsi come un mero strumento in questo processo. Di conseguenza, il proletariato è l'unica classe che ha un interesse costante e incondizionato nel portare a compimento il ruolo storico del capitale, fino alla sua piena realizzazione.Ciò implica l'eliminazione delle premesse che perpetuano la divisione sociale in classi.

L’acuto osservatore ha saggiamente evidenziato come la realizzazione di un ideale comunista possa essere illustrata con innumerevoli esempi. Si pensi, ad esempio, all’opera di uno scienziato, o di un gruppo di scienziati, che si adoperano senza sosta, giorno e notte, per offrire alla società il frutto della ricerca scientifica, dedicando, nel tempo libero, anche attività di divulgazione. In tali circostanze, la disponibilità di un minimo di risorse materiali rappresenta una condizione necessaria, e non il fine ultimo dell’impegno comunista. Analogamente, l’attività di uno scrittore, di un artista, di un musicista, che si prodigano affinché l’umanità possa osservare, esperire, comprendere più a fondo la realtà, ampliando così le potenzialità umane, e non perseguendo meri interessi di guadagno o di affermazione personale, può essere considerata un’espressione di tale ideale. Naturalmente, non tutte le attività riconducibili a questo ambito possono essere definite comuniste in senso stretto.Tuttavia, qualora i tratti distintivi di un’organizzazione sociale e spontanea dell’attività umana, delineati nel capitolo precedente, si manifestino concretamente, allora ci troveremmo al cospetto di un comunismo di scala ridotta, sebbene in questo specifico contesto esso assuma una forma piuttosto limitata.

Il comunismo, come giustamente si avrà modo di comprendere, è ben lungi dall'essere un progetto compiuto. In primo luogo, le premesse per la sua realizzazione si fondano su un mero e semplice sfruttamento capitalista dell'uomo sull'uomo. In secondo luogo, il capitale, in quanto rapporto sociale egemone, delimita anche in questo contesto la libertà entro schemi rigidi. Non è possibile edificare il comunismo, inteso come sistema di produzione materiale, senza superare i confini della sfera ideale. E sebbene il comunismo, nella sua dimensione ideale, abbia svolto un ruolo cruciale nel corso della storia umana – poiché è nell'attività libera che germinano le idee che non solo riflettono, ma anche "plasmano" la realtà – ciò non può comunque costituire un elemento determinante, poiché non trasforma ancora la società in una società comunista. In terzo luogo, si tratta qui di un superamento ancora parziale della divisione del lavoro a livello sociale, che rappresenta una condizione imprescindibile per l'avvento del comunismo.In questo contesto, l’universalità dell’uomo può esprimersi unicamente attraverso la coscienza, ma non si estende all’integrità dell’individuo inteso come entità pratica e concreta.

Esiste, tuttavia, una specifica modalità di produzione materiale che trascende la mera attività materiale, configurandosi come un’espressione completa della natura umana. Questa forma di produzione non si limita a influenzare indirettamente, ma comprende attivamente la genesi dell’ideale, della coscienza collettiva, del modello sociale, della definizione degli scopi comunitari. L’esempio precedentemente menzionato della giornata di lavoro collettiva, in cui gli operai si dedicavano al carico e scarico delle merci, nonché alla manutenzione del materiale ferroviario, ne offre un’efficace illustrazione. Il lavoro manuale, apparentemente di scarso fascino, non era riducibile a una semplice attività di movimentazione, bensì rappresentava un impegno volto alla costruzione di una rinnovata realtà sociale.Vale la pena ricordare le giornate dedicate alla raccolta della canna da zucchero a Cuba: muniti di strumenti modesti e operando sotto un sole implacabile, i partecipanti non si limitavano al semplice taglio della canna, ma raccoglievano anche i fondi necessari per ammodernare la produzione, erigere scuole e ospedali, al fine di creare le condizioni per il progresso di tutti i cittadini cubani. Questo genere di attività veniva praticato anche nella colonia intitolata a Gorkij e nella comune intitolata a Dzeržinskij, entrambe guidate da Anton Semënovič Makarenko. Ogni compito, persino lo scavo di latrine con una pala e la costruzione di piccole abitazioni in legno sopra di esse, si trasformava, in ultima analisi, in un processo di formazione di un nuovo individuo, in un’attività umana libera, ispirata ai principi di verità, bene e bellezza.

Di conseguenza, la questione della forma libera e direttamente collettiva dell’attività nella produzione materiale deve essere affrontata in questi termini: quale rapporto intrattiene tale attività con la società, intesa come totalità in divenire, e quale ruolo e importanza assume nel processo di sviluppo della totalità sociale?

Non è difficile scorgere come gli esempi finora citati affondino le radici nella storia del socialismo. In questo contesto, la mera formalizzazione della socializzazione dei mezzi di produzione elimina una serie di ostacoli che si frappongono a un libero e fluido mutamento della realtà sociale. Si potrebbe quindi obiettare che tali esempi non siano del tutto calzanti. Tuttavia, in primo luogo, l'eliminazione di queste barriere rende possibile la realizzazione di ciò che prima non esisteva, o per il quale le condizioni non erano ancora mature, ma piuttosto di ciò per il quale mancava l'insieme completo delle premesse necessarie, e l'unico elemento mancante era proprio l'abolizione del dominio del capitale. Il fatto che la rimozione di tali ostacoli apra la strada allo sviluppo delle tendenze comuniste implica, a ben vedere, che tali tendenze esistano già, e che il potenziale per un tale sviluppo sia in realtà presente all'interno del capitalismo stesso.

Si potrebbe obiettare che una giornata di lavoro volontario rappresenta un’esperienza effimera, un evento isolato, incapace di incidere in modo significativo sul corso della vita. Benché in tale contesto vengano perseguiti e realizzati sia obiettivi di natura sociale sia aspirazioni personali dei partecipanti, al di fuori di esso, nel tessuto della produzione materiale, continuano ad applicarsi le medesime leggi del capitalismo: la forza lavoro rimane una merce e ogni singolo proletario stringe un tacito patto di scambio con la propria classe, che si presenta come un capitale collettivo; la legge del valore mantiene la sua validità, a meno che la sua applicazione non venga esplicitamente limitata, esasperando così la tendenza intrinseca del capitale all’automazione dei processi produttivi.

In sostanza, l’esempio precedentemente menzionato, preso isolatamente, non può definirsi una società comunista. Le osservazioni formulate sono del tutto condivisibili e non intendiamo metterle in discussione. Tuttavia, l’aspetto fondamentale risiede non tanto nella mera constatazione del fatto che “questo è comunismo, quest’altro no”, o che “questo non è ancora comunismo”, quanto piuttosto nelle tendenze che animano la società e nel potenziale che si manifesta in queste piccole esperienze di comunismo. È con questa prospettiva che Lenin concepì l’idea dei “germogli del comunismo”, osservandoli proprio come tali, come embrioni. Per tale ragione, egli focalizzò la sua attenzione principalmente su due aspetti: 1) la libera espressione della volontà dei partecipanti alle giornate di lavoro volontario (i cosiddetti “subbotnik”), sia nei confronti dei mezzi di produzione sia nei confronti di se stessi; 2) il significativo incremento della produttività del lavoro durante queste giornate, un aumento intrinsecamente legato a questa libera manifestazione della volontà.

È doveroso sottolineare un ulteriore aspetto: negli esempi precedentemente citati, l'attività viene svolta con strumenti ben lontani dall'essere all'avanguardia, persino secondo i parametri del capitalismo, figuriamoci con una tecnologia di matrice comunista. Questa notevole mobilitazione di risorse umane è piuttosto il risultato di un arretrato sviluppo sociale: una carenza di tecnologie, una scarsa organizzazione dei processi produttivi, persino se considerati in una prospettiva strettamente capitalistica. A nostro avviso, questa rappresenta l'obiezione più pertinente contro un comunismo inteso come semplice prototipo di una nuova società. Il comunismo dettato dalla necessità, di cui si è discusso in precedenza, non costituisce ancora un comunismo scaturito da un'esigenza intrinseca e, di conseguenza, non implica di per sé lo sviluppo di una società comunista. Esso rappresenta, pur essendo una condizione rilevante, unicamente un preludio a tale società.Anche in questo caso, è imprescindibile esaminare il nesso e la continuità tra questa forma di comunismo e quella sviluppatasi su basi autonome, un tema che indubbiamente meriterebbe un approfondimento specifico.

Ciononostante, il richiamo alla necessità di un’analisi mirata di tale problematica non ci esime dal dovere di indicare, in termini generali, la concreta possibilità di una tale interconnessione, né dall’onere di fornire elementi a sostegno dell’esistenza di forme collettive e dirette di attività umana, fondate sulle più avanzate tecnologie del nostro tempo.

Il nesso tra il «comunismo dettato dalla necessità» e un comunismo fondato su principi autonomi, la loro continuità, può risiedere, innanzitutto, nell’educazione delle masse, nella progressiva acquisizione dell’abitudine a manifestare liberamente la propria volontà, orientandola verso le esigenze sociali che plasmano la vita collettiva; in secondo luogo, nella creazione di strutture comunitarie, capaci di organizzare la vita sociale, le quali, anche una volta superate le condizioni materiali che ne hanno determinato la nascita, non dovrebbero necessariamente dissolversi. Piuttosto, grazie all’educazione che in esse si concretizza, orientata ai principi di una libertà intesa come consapevolezza della necessità sociale, tali strutture potrebbero evolvere in forme di organizzazione comunitaria destinate a plasmare la vita sociale. Non coltiviamo utopie. Non ci prefiggiamo di delineare in dettaglio come ciò avverrà. Ogni volta, il processo dipenderà dalla complessa interazione delle circostanze storiche e concrete.Ci preme sottolineare la possibilità di un simile sviluppo futuro, che già ora si intravede in embrioni di questo genere all’interno delle strutture organizzative collettive dell’attività umana, e che implica una partecipazione libera e paritaria.

Riguardo al comunismo, inteso come sistema fondato su tecnologie avanzate, questi germi si manifestano nel processo di utilizzo condiviso delle moderne tecnologie dell’informazione e nella loro stessa creazione. A titolo d’esempio, il sistema operativo Linux e altri programmi sviluppati nell’ambito del movimento del software libero si fondano su principi comunisti.

Chiariremo subito che non ci riferiamo all’insieme dei software open source. Numerosi progetti moderni di grande portata hanno dato vita a consorzi di dimensioni considerevoli, e il fatto che il codice sia accessibile non implica che le attività di sviluppo siano esenti da vincoli. In tali circostanze, il consorzio esercita un controllo completo sulle direttive del progetto, ne finanzia le attività e ne beneficia in prima persona. Il codice è sviluppato da programmatori retribuiti. Non si tratta, quindi, di un modello di comunismo puro, sebbene sussista la possibilità di evolvere verso una dinamica simile. Ad esempio, possono essere realizzati prodotti di supporto che si rivelano utili per molteplici aziende contemporaneamente. Questa cooperazione consente di ridurre i costi di sviluppo. Tuttavia, in ultima analisi, il prodotto viene concepito con l’obiettivo di generare profitto. Spesso, per poter utilizzare un tale prodotto, e a maggior ragione per scopi commerciali, è necessario disporre di ingenti risorse, che sono solitamente accessibili solo alle grandi aziende.In genere, si tratta di un servizio o di un software destinato alla vendita, e ciò che viene sviluppato rappresenta una delle «componenti» essenziali, ovvero un elemento imprescindibile per il suo funzionamento e per il successo commerciale. Tuttavia, le grandi aziende hanno accaparrato le tecnologie che, sul piano tecnico, presentano un notevole potenziale, superando i confini del sistema capitalista. Questo potenziale risiede nella possibilità, concretizzatasi grazie allo sviluppo di macchine e reti informatiche, di creare e utilizzare software in modo esteso.

È nato il concetto di software libero, e attualmente esistono diversi progetti che ne incarnano i principi, operando in contesti privi di dinamiche commerciali. Inoltre, i programmatori si adoperano affinché nessuno ne ricavi vantaggi economici. In sostanza, si tratta di un esempio di comunismo applicato nella pratica, inserito all'interno del sistema capitalista. Il software open source garantisce un accesso totalmente libero ai programmi e al codice sorgente, offrendo al contempo la possibilità di distribuire liberamente sia i programmi che il codice sorgente, a condizione che siano stati creati utilizzando risorse open source già disponibili. Un esempio emblematico è rappresentato dal sistema operativo Linux, che si è dimostrato un valido concorrente rispetto al sistema operativo Windows, a pagamento. Le aziende cercano in ogni modo di esercitarne il controllo e sviluppano prodotti basati sul codice open source di questo sistema. Tuttavia, il sistema stesso rimane ad accesso gratuito.

È innegabile che il software libero e non profit, sovente, non sempre risponda alle aspettative, come evidenziato dagli utenti. Tuttavia, anche il software commerciale, prima di giungere alla sua versione definitiva, attraversa una fase di sviluppo simile, e talvolta alcuni progetti vengono interrotti per mancanza di prospettive concrete. La dinamica è la stessa, ma con una differenza cruciale: in questo caso, tutte le fasi del processo sono aperte e accessibili a chiunque. Anche ciò che appare incompleto o destinato al fallimento viene reso di dominio pubblico. Il software libero è, per sua natura, un prodotto in divenire, concepito per raggiungere la sua forma compiuta grazie al contributo degli utenti. Questo aspetto, inevitabilmente, ne limita la competitività nel contesto capitalistico, in un mondo caratterizzato da un livello di alfabetizzazione digitale ancora insufficiente.

Le aziende, nella progettazione dei loro prodotti, tengono conto del livello di alfabetizzazione dei potenziali utenti, mirando a creare soluzioni semplici e intuitive. Il risultato è un prodotto pensato per un contesto in cui, a fronte di una crescente alienazione, è più facile e immediato pagare per un’esperienza utente intuitiva, anziché investire tempo nell’apprendimento di un software più complesso. Questi prodotti non prevedono che l’utente crei o modifichi programmi esistenti, adattandoli alle proprie esigenze attraverso l’aggiunta o la riscrittura di codice. Il loro utilizzo non richiede un elevato grado di competenza, né tantomeno la comprensione dei processi tecnologici sottostanti. Di conseguenza, i programmi sviluppati dalle grandi aziende hanno maggiori probabilità di successo commerciale. Inoltre, spesso si distinguono per una qualità superiore, frutto di un’attenzione meticolosa ai dettagli.

Al contrario, il software libero e non commerciale richiede una conoscenza almeno elementare dei principi di programmazione, presupponendo la possibilità – e talvolta l’effettiva necessità – di aggiungere o modificare il codice sorgente.L’obiettivo è stimolare l’utente a sviluppare, almeno in parte, competenze di programmazione, incoraggiandone un continuo miglioramento in tale ambito\footnote{È indubbio che tale dinamica possa disorientare e ostacolare la concentrazione sul compito da svolgere. In ultima analisi, invece di dedicarsi compiutamente al proprio lavoro, gli individui si trovano costantemente impegnati a riparare gli strumenti che, inevitabilmente, si danneggiano nell’uso di Linux.Non esistono difficoltà insormontabili – almeno, se si considera il potenziale a disposizione – per risolvere questa situazione e creare programmi di qualità superiore, a patto che si possa investire più tempo nella loro elaborazione.}. Ed è proprio questo aspetto, che mitiga la spietata logica competitiva del software libero, a rivelarne il più profondo potenziale sociale rispetto al software commerciale (non si tratta di rinunciare a un’interfaccia intuitiva, bensì di consentire la creazione autonoma di un’interfaccia che lo sia). Nella sua creazione, nei suoi metodi di diffusione e nel suo utilizzo, risiede non solo il potenziale per un avanzamento tecnico, ma anche per uno sviluppo umano: superare la frammentazione del lavoro, elevando tutti i partecipanti al processo al livello di una competenza informatica moderna. E, sebbene questo embrionale modello di collaborazione venga spesso tentato di essere cooptato (e talvolta con successo) dalle aziende capitalistiche, esso può essere considerato uno dei primi germogli di una nuova società, ancora oggi fragile e in divenire. In primo luogo, esso è nato grazie all’impiego delle tecnologie più avanzate a disposizione del capitalismo.In secondo luogo, la traduzione definisce con precisione il confine tra l'impiego e lo sfruttamento capitalista di queste tecnologie e un loro utilizzo che trascende i limiti del sistema. È altresì fondamentale evidenziare come questa specifica modalità di produzione, di per sé, conduca all'abolizione delle gerarchie e alla convergenza tra utenti e sviluppatori. Essa non richiede collegamenti esterni tra gli individui, né la presenza di organi di controllo dedicati a garantirne la coesione. Il controllo dall'alto, anche nella sua forma più pura e democratica – ovvero attraverso la risoluzione delle problematiche mediante votazioni universali – non è più imprescindibile. In realtà, questa tendenza si sta affievolendo, divenendo progressivamente obsoleta.

È doveroso osservare che, quando si discute della produzione di software libero e non commerciale, ci si riferisce, ancora una volta, alla sfera della produzione ideale. In effetti, il processo di creazione del codice sorgente, inteso come una forma di testo, appartiene all'attività ideale, e il codice del programma (ma non il programma stesso, nella sua esecuzione ordinaria) costituisce un prodotto ideale. Tuttavia, si tratta di un'attività ideale che rappresenta una componente essenziale della moderna produzione industriale materiale. La creazione di software è un'attività finalizzata all'introduzione dell'automazione nel processo produttivo. E, benché il codice sorgente sia un prodotto ideale, il funzionamento delle moderne macchine industriali, basato su tali programmi (ossia il "funzionamento" dei programmi nel contesto industriale), è una produzione materiale e meccanica a tutti gli effetti.

È innegabile che, nell’ambito dello sviluppo di software libero, sorgano numerosi progetti superflui e privi di utilità. Tuttavia, questo non è l’aspetto cruciale. L’elemento significativo risiede nel fatto che, parallelamente a tali iniziative, vengono realizzati prodotti che possono effettivamente trovare un’applicazione su vasta scala nell’industria contemporanea. E questi prodotti nascono dalla collaborazione. Benché tali progetti non sempre giungano a compimento, o si protraggano per anni, ciò è spesso dovuto al fatto che i partecipanti devono provvedere al proprio sostentamento e, di conseguenza, dispongono di meno tempo da dedicare ad attività non remunerate.

Possiamo osservare uno schema analogo nel processo di creazione e utilizzo dei moderni database. Anche in questo ambito, non tutti i prodotti derivano da un approccio collettivista. Ciononostante, esistono vasti database di pubblico dominio che si fondano esclusivamente su una logica collaborativa, sia per quanto riguarda la loro creazione sia per il loro utilizzo. Un approccio diverso sarebbe impensabile, poiché ne comprometterebbe il valore e la fruibilità rispetto agli scopi per i quali sono stati concepiti. A titolo di esempio, alcuni database geospaziali incoraggiano la partecipazione libera alla loro elaborazione, con l'obiettivo di costruire una rappresentazione dello spazio utile per il progresso della società.

Anche in questo caso, le leggi oggettive che governano l'attività umana, nella loro intrinseca tendenza, escludono l'inevitabile gerarchia tipica del sistema di produzione capitalista, includendo tutti i partecipanti come soggetti uguali e liberi.

Si potrebbe obiettare che l’architettura del database rappresenti anch’essa un costrutto ideale. Tuttavia, nel processo di creazione di un database, l’attività concettuale è indissolubilmente legata alla sua realizzazione materiale – compresa la sua rappresentazione fisica all’interno del sistema – e alla riproduzione delle dinamiche sociali, tanto che, pur nella sua natura ideale, non può essere considerata un’attività appartenente alla sovrastruttura. Analogamente a quanto accade con il software, si tratta di un momento ideale della produzione sociale materiale, e non di un’astrazione – non nel senso di un’astrazione puramente concettuale, ma in un senso più ampio e reale – da tale produzione, ridotta a forme distinte della coscienza sociale.

Non ci riferiamo, dunque, a una mera aspirazione ideale, bensì alla concreta produzione industriale, alle tendenze e alle potenzialità in essa latenti fin dall’epoca del capitalismo. Queste, una volta superate le limitazioni imposte dalla proprietà privata dei mezzi di produzione, possono evolvere verso il comunismo, fornendo una solida base tecnologica e organizzativa, nonché un substrato educativo e formativo, per la nascita di un nuovo tipo di individuo.

L’ulteriore sviluppo di tale tendenza conduce alla fabbrica automatizzata, che, nell’ottica di Marx, rappresenta il prototipo della futura produzione materiale. In questo contesto si intravede anche la possibilità di realizzare le intuizioni di V. M. Glushkov in merito ai sistemi di gestione economica “a livello statale” (universale).

(questa fabbrica automatica concepita su scala sociale). Nel futuro (a patto che si superino gli ostacoli di ordine sociale), ogni individuo potrà accedere, in tempo reale, alle informazioni sullo stato dell’intera produzione sociale. Tutti potranno partecipare alla gestione non solo della produzione di beni – che ha senso unicamente in quanto strumento al servizio dello scopo primario della produzione stessa – ma anche della produzione dell’essere umano, della vita sociale nella sua interezza.

È interessante osservare come tale partecipazione a tutti questi processi implichi la creazione e la fruizione, da parte di ciascun individuo, di un prodotto ideale, un prodotto che, tuttavia, non sia alienato e che non possa, in alcun modo, essere alienato dalla produzione materiale. In ciò si può intravedere l’emergere di una nuova forma di coscienza sociale, che trascende le sue precedenti manifestazioni alienate.

È il capitalismo moderno a generare l’esigenza di superare gli ostacoli sociali, aprendo la strada a una libertà di azione. Esso stesso, al di là di ogni volontà o coscienza individuale, pone con forza questa questione al centro del dibattito. La logica intrinseca dell’evoluzione del capitale contemporaneo, che incessantemente alimenta conflitti, rende impraticabili le modalità e le forme consolidate di attività umana. Il capitale spinge, quindi, verso la creazione di nuove e più efficaci modalità operative, volte a risolvere le problematiche connesse alla riproduzione della vita sociale, opponendosi alla sua intrinseca forza distruttiva
\end{document}
